\chapter{Distributions and Frobenius' theorem}

\begin{definition}[\emph{distributions}]
    Let \(M\) be a smooth manifold.
    A \(k\)-dimensional \emph{distribution} \(\cal{V}\) is a map \(M\to\{\text{subspaces of }T_pM\}\) which assigns each \(p\in M\) to a \(k\)-dimensional subspace \(\cal{V}_p\subseteq T_p M\).
\end{definition}

\begin{definition}[\emph{smooth distribution}]
    A distribution \(\cal{V}\) is said to be \emph{smooth} if for each \(p\in M\), there exists a neighborhood \(U\) of \(p\) and smooth vector fields \(V_1,\dots,V_k\) on \(U\) such that \(\cal{V}_q = \spn\{V_1(q),\dots,V_k(q)\}\) for all \(q\in U\).
\end{definition}

\begin{definition}
    A vector field \(v\) \emph{belongs to a distribution} \(\cal{V}\) if \(v(p)\in\cal{V}_p\) for all \(p\in M\).
\end{definition}

\begin{definition}[\emph{integral manifold1}]
    A submanifold \(N\subset M\) is called an \emph{integral manifold} of \(\cal{V}\) if for every \(p\in M\), \[
        \im(d\iota_N:T_pN\to T_pM) = \cal{V}_p.
    \]
\end{definition}

\begin{example}
    Let \(V_1 = \spn\{(1,0,0), (0,1,0)\}\) be a 2-\(\dim\) distribution on \(\R^3\).
    \[
        \begin{tikzpicture}[
    >=Stealth,
    axis/.style={
        ->,
        black,
        line width=1pt
    },
    xone/.style={
        ->,
        red!85!black,
        line width=1pt
    },
    xtwo/.style={
        ->,
        blue!80!black,
        line width=1pt
    },
    gridpt/.style={
        fill=black,
        inner sep=1pt,
        circle
    }
]

% ------------------------------------------------
% Axes in R^3
% ------------------------------------------------
\coordinate (O) at (0,0);

% x-axis
\coordinate (ex) at (-2.0,-1.5);

% y-axis
\coordinate (ey) at (2.6,0);

% z-axis
\coordinate (ez) at (0,2.3);

\draw[axis] (O) -- (ex) node[below left] {\(x\)};
\draw[axis] (O) -- (ey) node[right] {\(y\)};
\draw[axis] (O) -- (ez) node[above] {\(z\)};

\node[below right] at (O) {\(0\)};

% ------------------------------------------------
% Vector directions parallel to x-axis and y-axis
% ------------------------------------------------
\coordinate (Xdir) at (-0.45,-0.3375); % parallel to x-axis
\coordinate (Ydir) at (0.55,0);        % parallel to y-axis

% ------------------------------------------------
% Points in the xy-plane where the distribution is shown
% ------------------------------------------------
\foreach \p in {
    % (-1.25,-0.85),
    % (-0.55,-0.85),
    % (0.15,-0.85),
    % (0.85,-0.85),
    (-1.25,-0.35),
    (-0.55,-0.35),
    (0.15,-0.35),
    % (1.25,-0.35),
    (-0.65,0.20),
    (0.05,0.20),
    (0.75,0.20),
    (-0.05,0.75),
    (0.65,0.75),
    (1.35,0.75)
} {
    \fill[gridpt] \p circle (0.7pt);

    % Red arrows parallel to x-axis
    \draw[xone] \p -- ++(Xdir);

    % Blue arrows parallel to y-axis
    \draw[xtwo] \p -- ++(Ydir);
}

% ------------------------------------------------
% Labels for basis vectors
% ------------------------------------------------
% \draw[xone] (2.55,1.35) -- ++(Xdir)
%     node[midway, above left] {\(X_1\)};

% \draw[xtwo] (2.55,0.85) -- ++(Ydir)
%     node[midway, above] {\(X_2\)};

% \node[right] at (3.05,1.08) {
%     \(\operatorname{span}\{X_1,X_2\}\)
% };

\end{tikzpicture}
    \]
\end{example}




\begin{lemma}
  Given \(F:\R\times\R\to\R\), let \(g(t) = F(t,t)\). Then \[
    \frac{dg}{dt}(t_0) = \frac{\pa F}{\pa s_1}(s_1,t_0)\bigg|_{s_1 = t_0} + \frac{\pa F}{\pa s_2}(t_0,s_2)\bigg|_{s_2 = t_0}.
  \]
\end{lemma}

\begin{theorem}
  Let \(v\) be a complete vector field on a smooth manifold \(M\), then \begin{enumerate}
    \item \(\L_v f = D_v f\) for any function \(f\);
    \item \(\L_v w = [v,w]\) for any vector field \(v\).
  \end{enumerate}
\end{theorem}
\begin{proof}
  \begin{enumerate}
    \item Let \(\c_p(t)\) be an integral curve of \(v\) with \(\c_p(0) = v\), then \[
      \phi_t^*f = f(\phi_t(p)) = f(\c_p(t)).
    \]
    Now \[
      \L_v f(p) = \frac{d}{dt}\phi_t^* f(p) = \frac{d}{dt}f(\c_p(t)) = D_v f(p).
    \]
    \item We have \[
      D_{\L_v w}f = D_{[v,w]}f = D_v D_w f - D_w D_v f.
    \]
    Since \(M\xrightarrow{\phi_{-t}}M\) is a diffeomorphism, \(\phi_{-t}\)-related implies \[
      D_{(\phi_{-t})_*w} = \phi_t^*D_w\phi_{-t}^*.
    \]
    Now \begin{align*}
      D_{\L_v w}f = \frac{d}{dt}D_{(\phi_{-t})_*w} f\bigg|_{t=0} 
      &= \frac{d}{dt}\underbrace{\phi_t^*D_w\phi_{-t}^*}_{F(t,t)}f\bigg|_{t=0} \\
      &= \frac{d}{dt}\phi_t^*D_w f\bigg|_{t=0} + \frac{d}{dt}D_w\phi_{-t}^*f\bigg|_{t=0} \\
      &= D_v D_w f - D_w D_v f = D_{[v,w]}f.
    \end{align*}
    Note that \begin{gather*}
      F(s_1,s_2) = \phi_{s_1}^*D_w\phi_{-s_2}^* \\
      F(s_1,0) = \phi_{s_1}^*D_w \\
      F(0,s_2) = D_w\phi_{-s_2}
    \end{gather*}
  \end{enumerate}
\end{proof}

\newpage
\begin{theorem}
  Suoppose \(v\), \(w\) are complete vector fields on a smooth manifold \(M\). Then \[
    \exp(tv) \text{ and } \exp(sw) \text{ commute } \forall s,t
    \iff [v,w] = 0.
  \]
\end{theorem}
\begin{proof}
  Recall that \[
    [v,w] = \L_v w = \frac{d}{dt} \underbrace{(\exp(-t))_* w}_{w_t} \bigg|_{t=0}.
  \]
  A common trick is \[
    \frac{d}{dt}w_t\bigg|_{{t=a}}
    = \frac{d}{dt} \exp(-av)\exp(-sv)
  \]
  Here \(
    w_t = (\exp(-tv))_* w
  \) is independent of \(t\). At \(t = 0\), \[
    \exp(-tv)_* = \id \implies w_t = w_0 = w \forall t.
  \]
\end{proof}

\begin{lemma}[Functoriality]
  If we have a smooth map \begin{align*}
    M\xrightarrow{\f}N \\
    v \leadsto w
  \end{align*} where \(v,w\) are complete. Then \(v,w\) are \(\f\)-related if and only if \[
    \begin{tikzcd}
      M \ar{r}{\f} \ar[']{d}{\exp(tv)} & N \ar{d}{\exp(tw)} \\
      M \ar{r}{\f} & N
    \end{tikzcd} \text{ commutes},
  \] i.e.\[
    \exp(tw)\circ\f = \f\circ\exp(tv)
  \]
\end{lemma}
\begin{proof}
  Let \(p\in M\), \(q = \f(p)\). Let \(\c_p:\R\to M\) be an integral curve of \(v\) with \(\c_p(0) = p\) and \(\c_q:\R\to N\) be an integral curve of \(w\) with \(\c_q(0) = q\). 
  We proved \(v,w\) are \(\f\)-related \[
    \iff \f(\c_p(t)) = \c_q(t) \quad\forall p\in M.
  \]
  By definition, \begin{align*}
    \c_p(t) &= \exp(tv)(p) \\
    \c_q(t) &= \exp(tw)(q).
  \end{align*}
  Hence, \begin{align*}
    \iff \f\circ\exp(tv)(p) = \exp(tw)\circ\f(p) \quad\forall p\in M.
  \end{align*}
\end{proof}

In particular, if we apply this lemma on the diffeomorphism \(\f:M\to M\).
\(\f\) preserves \(v\), i.e., \(\f_* v = v\) if and only if \(\f\) commutes with \(\exp(tv)\) for all \(t\).

\begin{theorem}
  Suppose \(v_1,v_2\) are complete vector fields such that \([v_1,v_2] = 0\). Then \begin{align*}
    \c:\R^2 &\to \operatorname{Diff}(M) \\
    (t_1,t_2) &\mapsto \exp(t_1v_1)\circ\exp(t_2v_2)
  \end{align*} is a (two-parameter) group homomorphism, i.e. \[
    \c(t_1 + t_1', t_2 + t_2') = \c(t_1,t_2)\circ\c(t_1',t_2').
  \]
\end{theorem}
The above is a very satisfactory answer to when flows commute.

We are going to discuss a more interesting part of our course. Let's talk about applications of the Lie bracket.

Intuitively, these are just higher dimensional vector fields.

\begin{definition}[\emph{distributions}]
    Let \(M\) be a smooth manifold.
    A \(k\)-dimensional \emph{distribution} \(\cal{V}\) is a map \(M\to\{\text{subspaces of }T_p M\}\) which assigns each \(p\in M\) to a \(k\)-dimensional subspace \(\cal{V}_p\subseteq T_p M\).
\end{definition}

\begin{definition}[\emph{smooth distribution}]
    A distribution \(\cal{V}\) is said to be \emph{smooth} if for each \(p\in M\), there exists a neighborhood \(U\) of \(p\) and smooth vector fields \(V_1,\dots,V_k\) on \(U\) such that \(\cal{V}_q = \spn\{V_1(q),\dots,V_k(q)\}\) for all \(q\in U\).
\end{definition}

\begin{definition}
    A vector field \(v\) \emph{belongs to a distribution} \(\cal{V}\) if \(v(p)\in\cal{V}_p\) for all \(p\in M\).
\end{definition}

\begin{definition}[\emph{integral manifold1}]
    A submanifold \(N\subset M\) is called an \emph{integral manifold} of \(\cal{V}\) if for every \(p\in N\), \[
        \im(d\iota_N:T_p N\to T_p M) = \cal{V}_p.
    \]
\end{definition}

\begin{example}
    Let \(\cal{V}_1\) be a 2-dimensional distribution on \(\R^3\) generated by \[
      X_1 = \<1,0,0\>, \quad X_2 = \<0,1,0\>.
    \]
    \[
        \begin{tikzpicture}[
    >=Stealth,
    axis/.style={
        ->,
        black,
        line width=1pt
    },
    xone/.style={
        ->,
        red!85!black,
        line width=1pt
    },
    xtwo/.style={
        ->,
        blue!80!black,
        line width=1pt
    },
    gridpt/.style={
        fill=black,
        inner sep=1pt,
        circle
    }
]

% ------------------------------------------------
% Axes in R^3
% ------------------------------------------------
\coordinate (O) at (0,0);

% x-axis
\coordinate (ex) at (-2.0,-1.5);

% y-axis
\coordinate (ey) at (2.6,0);

% z-axis
\coordinate (ez) at (0,2.3);

\draw[axis] (O) -- (ex) node[below left] {\(x\)};
\draw[axis] (O) -- (ey) node[right] {\(y\)};
\draw[axis] (O) -- (ez) node[above] {\(z\)};

\node[below right] at (O) {\(0\)};

% ------------------------------------------------
% Vector directions parallel to x-axis and y-axis
% ------------------------------------------------
\coordinate (Xdir) at (-0.45,-0.3375); % parallel to x-axis
\coordinate (Ydir) at (0.55,0);        % parallel to y-axis

% ------------------------------------------------
% Points in the xy-plane where the distribution is shown
% ------------------------------------------------
\foreach \p in {
    % (-1.25,-0.85),
    % (-0.55,-0.85),
    % (0.15,-0.85),
    % (0.85,-0.85),
    (-1.25,-0.35),
    (-0.55,-0.35),
    (0.15,-0.35),
    % (1.25,-0.35),
    (-0.65,0.20),
    (0.05,0.20),
    (0.75,0.20),
    (-0.05,0.75),
    (0.65,0.75),
    (1.35,0.75)
} {
    \fill[gridpt] \p circle (0.7pt);

    % Red arrows parallel to x-axis
    \draw[xone] \p -- ++(Xdir);

    % Blue arrows parallel to y-axis
    \draw[xtwo] \p -- ++(Ydir);
}

% ------------------------------------------------
% Labels for basis vectors
% ------------------------------------------------
% \draw[xone] (2.55,1.35) -- ++(Xdir)
%     node[midway, above left] {\(X_1\)};

% \draw[xtwo] (2.55,0.85) -- ++(Ydir)
%     node[midway, above] {\(X_2\)};

% \node[right] at (3.05,1.08) {
%     \(\operatorname{span}\{X_1,X_2\}\)
% };

\end{tikzpicture}
    \]
    This is just a plane parallel to the \(xy\)-plane. The integral manifolds of \(\cal{V}_1\) are just planes parallel to the \(xy\)-plane.
\end{example}

\begin{example}
    Let \(\cal{V}_2\) be a 2-dimensional distribution on \(\R^3\) generated by \[
      X_1 = \<-y,x,0\>, \quad X_2 = \<0,0,1\>.
    \]
    This is a circle in the \(xy\)-plane and a line parallel to the \(z\)-axis. The integral manifolds of \(\cal{V}_2\) are just cylinders parallel to the \(z\)-axis.
\end{example}

\begin{example}
  Let \(\cal{V}_3\) be a 2-dimensional distribution on \(\R^3\) generated by \[
    X_1 = \<1,0,y\>, \quad X_2 = \<0,1,0\>.
  \]
  There is no integral manifold through origin.
\end{example}
\begin{proof}
  Suppose \(N\) is an integral submanifold through origin, then it must also contain the integral curve of \(X_1\), which is a piece of \(x\)-axis. Similarly, it must also contain the integral curve of \(X_2\) passing \((t,0,0)\), which contains a piece of line segment parallel to the \(y\)-axis. Hence, \(N\) must contain a piece of the \(xy\)-plane containing the origin, this is a contradiction as \(\<1,0,0\>\not\in\cal{V}_p\) when \(p\) is not on the \(x\)-axis.
\end{proof}

We want to know when a distribution has an integral manifold. This is the content of Frobenius theorem.

The following is an algebraic condition that is easy to check.
\begin{definition}[\emph{involutive distribution}]
    A smooth distribution \(\cal{V}\) is said to be \emph{involutive} if it satisfies the \emph{Frobenius condition}, i.e., if \(X,Y\) are vector fields belong to \(\mathcal{V}\), so is \([X,Y]\).
\end{definition}

Note that for all \(f,g\in C^\infty(M)\), we write \[
  [fX, gY] = (fg)[X,Y] + (f\cdot D_X g)Y - (g\cdot D_Y f)X,
\] since \([X,Y]\), \(Y\) and \(X\) belong to \(\cal{V}\), we only need to check the condition on a set of basis.

\begin{example}
  For \(\cal{V}_1\), we have \([X_1,X_2] = 0\).
\end{example}
\begin{example}
  For \(\cal{V}_2\), we have \begin{align*}
    D_{X_1} &= -y\frac{\pa}{\pa x} + x\frac{\pa}{\pa y} \\
    D_{X_2} &= \frac{\pa}{\pa z}
  \end{align*}
  Then \[
    D_{[X_1,X_2]} = D_{X_1}D_{X_2} - D_{X_2}D_{X_1} = 0 \implies [X_1,X_2] = 0,
  \] so \([X_1,X_2] = 0\).
\end{example}

\begin{example}
  Cant see.
\end{example}

\begin{theorem}
  If \(\cal{V}\) is integrable, then it is involutive.
\end{theorem}
If you have a submanifold \(N\) in the manifold \(M\). Given two vector fields, you can define the Lie bracket using the vector fields on the submanifold, hence involutive.
\begin{proof}
  Suppose \(X\) and \(Y\) are vector fields belong to \(\cal{V}\). Then \(X\restriction_N\) and \(Y\restriction_N\) are vector fields tangent to \(N\) by the definition of integral manifold. Consider the \(\iota:N\to M\), \(X\restriction_N\) and \(X\) are \(\iota\)-related, similarly \(Y\restriction_N\) and \(Y\) are \(\iota\)-related.These imply \([X\restriction_N,Y\restriction_N]\) and \([X,Y]\) are \(\iota\)-related, i.e. \[
    [X,Y]_p \text{ and } 
  \]
\end{proof}